\subsection{从李代数态射到李群态射的过渡}

引理 1. 设 G 是单连通†拓扑群，W 是 e 的对称连通开邻域，H 是一个群，f 是从 W³ 到 H 的映射，使得

\[
f (x y z) = f (x) f (y) f (z)
\]

当 \(x, y, z\) 属于 W 时，存在一个态射 \(f'\)，它从 G 到 H，且满足 \(f' \mid \mathrm{W} = f \mid \mathrm{W}\)。对于 \((g, h) \in \mathrm{G} \times \mathrm{H}\) 以及 W 中 \(e\) 的一个开邻域 U，令 \(\mathrm{A}(g, h, \mathrm{U})\) 为由 \((gu, hf(u)) \in \mathrm{G} \times \mathrm{H}\)（其中 \(u \in \mathrm{U}\)）组成的集合。于是 \((g, h) \in \mathrm{A}(g, h, \mathrm{U})\)，且 \(\mathrm{A}(g, h, \mathrm{U}_1) \cap \mathrm{A}(g, h, \mathrm{U}_2) = \mathrm{A}(g, h, \mathrm{U}_1 \cap \mathrm{U}_2)\)。令 \((s, t) \in \mathrm{A}(g, h, \mathrm{U})\)；有 \(s = gu\) 且 \(t = hf(u)\)，其中 \(u \in \mathrm{U}\)；存在一个开邻域 \(\mathrm{U}'\)，它是 W 中 \(e\) 的邻域，使得 \(u\mathrm{U}' \in \mathrm{U}\)；于是对 \(u' \in \mathrm{U}'\)，

\[
(s u ^ {\prime}, t f (u ^ {\prime})) = (g u u ^ {\prime}, h f (u u ^ {\prime})) \in \mathrm{A} (g, h, \mathrm{U})
\]

从而 \(\mathrm{A}(s,t,\mathrm{U}^{\prime})\subset \mathrm{A}(g,h,\mathrm{U})\)。于是 \(\mathrm{A}(g,h,\mathrm{U})\) 构成 \(\mathbf{G}\times \mathbf{H}\) 上拓扑的一个基。我们把 \(\mathbf{Y}\) 记为带有此拓扑的 \(\mathbf{G}\times \mathbf{H}\)，并把规范投影 \(p\) 记为从 \(\mathbf{Y}\) 到 \(\mathbf{G}\) 的投影，且它是开映射。\(p\) 在 \(\mathrm{A}(g,h,\mathrm{U})\) 上的限制是从 \(\mathrm{A}(g,h,\mathrm{U})\) 到 \(g\mathrm{U}\) 的同胚。因此 \((\mathrm{Y},p)\) 是 \(\mathbf{G}\) 的覆盖空间。令 \(\mathrm{Y}_0\) 为 \(\mathrm{Y}\) 中由 \(\mathrm{A}(e,e,\mathrm{W})\) 生成的子群，令 \(\mathfrak{B}\) 为 \(\mathrm{A}(e,e,\mathrm{U})\) 的集合。显然，\(\mathfrak{B}\) 满足 General Topology, Chapter III, § 1, no. 2 的条件 \((\mathrm{GV}_{\mathrm{I}}^{\prime})\) 和 \((\mathrm{GV}_{\mathrm{II}}^{\prime})\)。集合 \(\mathrm{Y}_0^{\prime}\) 由满足下述条件的 \(y\in \mathrm{Y}_0\) 组成：映射 \(z\mapsto yzy^{-1}\) 和 \(z\mapsto y^{-1}zy\) 从 \(\mathrm{Y}_0\) 到 \(\mathrm{Y}_0\)，并在 \((e,e)\) 处连续；它是 \(\mathrm{Y}_0\) 的一个子群。令 \(w\in W\)。映射 \(w'\mapsto ww'w^{-1}\) 从 \(\mathrm{W}\) 到 \(\mathrm{G}\) 连续，因而映射

\[
\left(w ^ {\prime}, f \left(w ^ {\prime}\right)\right) \mapsto \left(w w ^ {\prime} w ^ {- 1}, f \left(w w ^ {\prime} w ^ {- 1}\right)\right)
\]

从 \(\mathrm{A}(e,e,\mathrm{W})\) 到 \(\mathrm{Y}_0\) 在 \((e,e)\) 处连续。现在

\[
f \left(w w ^ {\prime} w ^ {- 1}\right) = f (w) f \left(w ^ {\prime}\right) f \left(w ^ {- 1}\right)
\]

\[
\left(w w ^ {\prime} w ^ {- 1}, f \left(w w ^ {\prime} w ^ {- 1}\right)\right) = (w, f (w)) \left(w ^ {\prime}, f \left(w ^ {\prime}\right)\right) \left(w, f (w)\right) ^ {- 1}
\]

由于 \(w^{-1} \in W\)，可见 \((w, f(w)) \in Y_0'\)。因此 \(A(e, e, W) \subset Y_0'\)，从而 \(Y_0' = Y_0\)。于是群 \(Y_0\) 带有滤子基 \(\mathfrak{B}\)，并满足 General Topology, Chapter III, § 1, no. 2 的条件 \((GV_{\mathrm{III}}')\)。由于

\[
(g, h). \mathrm{A} (e, e, \mathrm{U}) = \mathrm{A} (g, h, \mathrm{U}),
\]

\(Y_{0}\) 是一个拓扑群，并且由于 \(A(e, e, W)\) 连通而连通。于是 \(p(Y_{0})\) 是 G 的开子群；由于 G 连通，故 \(p(Y_{0}) = G\)。\(p \mid Y_{0}\) 的核是离散的。由于 G 单连通，\(p \mid Y_{0}\) 是从 \(Y_{0}\) 到 G 的同胚。因此，\(Y_{0}\) 是从 G 到 H 的态射 \(f'\) 的图。对于 \(g \in W\)，\((g, f(g)) \in A(e, e, W) \subset Y_{0}\)，从而 \(f(g) = f'(g)\)。

定理 1. 设 G 和 H 是李群，h 是从 L(G) 到 L(H) 的连续态射。假设 G 单连通。则存在唯一一个从 G 到 H 的李群态射 \(\phi\)，使得 \(h = \mathrm{L}(\phi)\)。

\(\phi\) 的存在性由引理 1 和 §4, no. 1, Theorem 1 (i) 得出。\(\phi\) 的唯一性由 §4, no. 1, Theorem 1 (ii) 以及 \(G\) 连通这一事实得出。

推论。设 G 是有限维单连通李群。则存在一个有限维解析线性表示，其核是离散的。

根据 Chapter I, § 7, Theorem 2，存在有限维向量空间 E 以及一个单射态射 \(h\)，它从 \(L(G)\) 到李代数 \(\operatorname{End}(E)\)。由定理 1，存在一个态射 \(\phi\)，它从 G 到 \(\mathbf{GL}(E)\)，使得 \(L(\phi) = h\)。因此 \(\phi\) 是浸入，从而其核是离散的。

注记。(1) 存在没有有限维单射解析线性表示的有限维单连通李群（练习 2）。

\begin{enumerate}
\def\labelenumi{(\arabic{enumi})}
\setcounter{enumi}{1}
\tightlist
\item
  存在有限维连通李群 G，使 G 的每个有限维解析线性表示都有非离散的核（练习 3 和 4）。
\end{enumerate}

